\chapter*{Preface}
\addcontentsline{toc}{chapter}{Preface}

This book was written from building a thing: a notebook, in the spirit of Mathematica, that
differentiates, integrates, row-reduces and simplifies, and that \emph{shows its work} --- every
answer arrives with the chain of rewrites that produced it, every subterm of every line can be
clicked to ask where it came from, and every rule that fired carries a label saying how much of
it is a theorem. The engine behind the notebook is a single Lean~4 program. It runs natively as a
server and a command-line tool, and it runs in the browser as WebAssembly; the notebook never
computes anything itself. The proofs about the engine are Lean~4 too, and they are checked every
time the project builds.

The order of the chapters is roughly the order in which the thing was built, because the order
turned out to matter. A rewriter came first, and with it a trace. Then meaning: what each rule
claims, over the integers, then over the reals, then under a derivative. Then termination,
which was the hard part and is the centre of the book: the engine has no step budget, and the
proof that it needs none took several wrong orderings and one wall that would not move. Then the
tracing of origins, then three more ``worlds'' --- linear algebra, integration, the
$\lambda$-calculus and finite order theory --- each of which had to find its own honest place
between \emph{proved} and \emph{checked}.

\medskip
\textbf{Who this book is for.} Someone who can read undergraduate mathematics and some functional
code, who has perhaps met Lean~4 in one of its companion volumes, and who wants to see what
``verified'' can mean for a system that is also meant to be \emph{used}. No Mathlib is needed to
read the engine; where the reals enter, the book says so.

\medskip
\textbf{What is claimed.} Precisely what the notebook's status labels claim, and no more. A rule is
\status{verified} when an unconditional theorem says its output means what its input meant;
\status{conditional} when the theorem needs a hypothesis and the \emph{necessity} of the hypothesis
is itself proved; \status{checked} when the step is a guess that a later step verifies; and
\status{unverified} when there is no theorem yet. Chapter~\ref{ch:ledger} is the full ledger.
Several places where a proof was avoided are pointed at by name: they are the run-time checks the
book calls \emph{honest boundaries}, and the reader will see exactly what is checked and why a
check was the right thing there.

\medskip
\textbf{How to read this book.} Part~I is the engine's shape and the rewriter that keeps a trace.
Part~II gives the rules their meaning. Part~III is termination. Part~IV is the notebook side of
``show your work''. Part~V adds the other worlds. Boxes labelled \emph{From the log} retell what
actually happened --- the wrong turns are kept, since they are the argument for the right ones ---
and boxes labelled \emph{Try it} give cells to type into the notebook.

\medskip
\textbf{The code.} Everything quoted is from the repository at the commit stamped on the title
page. File names are given the way the repository has them: \file{engine/MathEngine/Order.lean}
is the ordering, \file{proofs/Proofs/SimpReal.lean} the real-number theorems, and so on.
